\section{A conditional all-order noncentral band}
The numerical growth of the local radii suggests a more
structural criterion, independent of individual moment
calculations. It is presented as a conditional theorem,
\emph{not} as an established property of the actual theta
kernel.

\begin{proposition}[Radially increasing potential curvature]\label{prop:radial}
Let $K(u)=e^{-V(u)}$ be positive and even, with sufficient
decay to have all moments. Suppose $V$ is twice continuously
differentiable and $V''(u)$ is nondecreasing for $u\ge0$.
Then the ratios $B_n/A_n$ of the preceding section
are nonincreasing as $n$ increases:
\[
\frac{B_{n+1}}{A_{n+1}}\le
\frac{B_n}{A_n},\qquad n\ge0.
\]
Consequently the local radii $\sqrt{2A_n/B_n}$
are nondecreasing, and the single order-one bound,
if valid, gives $L_n[F](x)>0$ for every $n\ge1$
throughout $|x|<\sqrt{2A_1/B_1}$.
\end{proposition}
\begin{proof}
Set $d=u-v$, $w=u+v$. Up to a constant, the
joint density on the positive quadrant $d,w\ge0$ is
\[
g(d,w)=\exp\{-V((w+d)/2)-V((w-d)/2)\}.
\]
By evenness, $V''$ is even and therefore
\[
\partial_d\partial_w\log g
=\frac14\left[V''\!\left(\frac{w-d}{2}\right)
                  -V''\!\left(\frac{w+d}{2}\right)\right]
\le0,
\]
because $|w-d|\le w+d$. Thus for $d_2>d_1$
the likelihood ratio $g(d_2,w)/g(d_1,w)$
is nonincreasing in $w\ge0$. By the elementary
reverse-Chebyshev covariance inequality, the
conditional expectation $h(d)=
\mathbb E[w^2\mid |u-v|=d]$ is nonincreasing in $d$.

Under the tilted joint density proportional to
$|d|^{2n}g(|d|,|w|)$, one has
$B_n/A_n=\mathbb E_n[w^2]=\mathbb E_n[h(d)]$.
Passing from $n$ to $n+1$ multiplies the marginal
density of $d$ by the increasing weight $d^2$,
so this expectation of the decreasing function
$h$ cannot increase. This proves the ratio
inequality; the noncentral band follows from
\eqref{eq:local-positive}.
\end{proof}

\paragraph{Unproved theta-kernel gate.}
For the genuine Jacobi theta kernel, the required
global hypothesis is
\[
V'''(u)=-\big(\log\Phi(u)\big)'''\ge0
\qquad(u>0).
\]
High-precision numerical tests at
$u=0.0001,0.001,0.01,0.025,0.05,0.1,0.2,
0.4,0.8,1.5,2$ returned strictly positive values;
this is not a rigorous interval or an all-$u$
certificate. If proved, it would yield a
nontrivial all-order Laguerre positivity \emph{strip}
for the actual transform, but still not RH.
Checking prior art and finding a falsifying region
for the curvature condition are required before
making a claim of novelty.
